%% LyX 2.4.4 created this file.  For more info, see https://www.lyx.org/.
%% Do not edit unless you really know what you are doing.
\documentclass[english]{article}
\usepackage[T1]{fontenc}
\usepackage[latin9]{inputenc}
\usepackage{enumitem}
\usepackage{amsmath}
\usepackage{amsthm}
\usepackage{amssymb}
\usepackage{cancel}
\usepackage{geometry}
\geometry{verbose,tmargin=1in,bmargin=1in,lmargin=1in,rmargin=1in}
\PassOptionsToPackage{normalem}{ulem}
\usepackage{ulem}

\makeatletter
%%%%%%%%%%%%%%%%%%%%%%%%%%%%%% Textclass specific LaTeX commands.
\numberwithin{equation}{section}
\numberwithin{figure}{section}
\newlength{\lyxlabelwidth}      % auxiliary length 

%%%%%%%%%%%%%%%%%%%%%%%%%%%%%% User specified LaTeX commands.
\usepackage{fancyhdr}
\usepackage{lastpage}
\pagestyle{fancy}
\usepackage{enumitem}
\usepackage{ifthen}
\everymath=\expandafter{\the\everymath\displaystyle}
%\DeclareMathSizes{10}{10}{10}{10}
\usepackage{multicol}

\usepackage{tikz}
\usetikzlibrary{patterns}
\usetikzlibrary{plotmarks}
\usepackage{pgfplots}
\pgfplotsset{compat=newest}

\usepackage{advdate}    % Advancing/saving dates
\usepackage{datetime}   % Dates formatting
\usepackage{datenumber} % Counters for dates
\usdate
% Added by lyx2lyx
\setlength{\parskip}{\smallskipamount}
\setlength{\parindent}{0pt}

\makeatother

\usepackage{babel}
\begin{document}
\renewcommand{\author}{Dr.\,James Ashe}

\newcommand{\classNum}{335}

\newcommand{\classSec}{A}

\newcommand{\examType}{Midterm Exam}

\renewcommand{\date}{\formatdate{4}{3}{2014}}

%%First page's header%%%%%%%%%%%%%%%%%%%%%%
\fancypagestyle{firststyle} %%header settings for first pagr
{
	\fancyhead[L]{MTH \classNum{}(\classSec{}) \author{}\\
	\textbf{\examType{}} \date{}}
	\fancyhead[C]{\textbf{Name:}}
	\setlength{\headsep}{14pt}
}
%%%%%%%%%%%%%%%%%%%%%%%%%%%%%%%%%%%%%%%%%%

%%Next page's header%%%%%%%%%%%%%%%%%%%%%%%
\setlist{nolistsep}
\lhead{\classNum{}(\classSec{})} 
\chead{\textbf{}} 
\cfoot{\thepage{}~of~\pageref{LastPage}} 
\setlength{\headsep}{5pt} 
%%%%%%%%%%%%%%%%%%%%%%%%%%%%%%%%%%%%%%%%%%%

%%Various settings; see the preamble for other settings%%%%%
%\setenumerate[0]{leftmargin=12pt,itemindent=0pt,labelwidth=0pt}
\setlist{topsep=.5pt}
\renewcommand{\labelenumi}{\textbf{\arabic{enumi}.}}
\renewcommand{\labelenumii}{\textbf{\alph{enumii}.}}
\thispagestyle{firststyle}
%%%%%%%%%%%%%%%%%%%%%%%%%%%%%%%%%%%%%%%%%%

\newcommand{\Dspace}{.8cm}

\emph{Choose any $12$ problems.}
\begin{enumerate}[resume]
\item Determine the following:

\begin{enumerate}
\item Write the following set in tabular form: $A=\{x\colon x\mbox{ is a letter in "do not repeat"}\}$.\vspace{\Dspace}
\item Write the following set in tabular form: $B=\{x\colon x=2n\mbox{ for }n\in\mathbb{Z}\}$.\vspace{\Dspace}
\item Write the set $C=\left\{ \dots,-3,0,3,6,\dots\right\} $ in set-builder
form.\vspace{\Dspace}
\end{enumerate}
\item Give the definition of the following:

\begin{enumerate}[resume]
\item A \emph{function.}\vspace{\Dspace}\vspace{\Dspace}
\end{enumerate}
\begin{enumerate}
\item A \emph{1-1 function}.\vspace{\Dspace}\vspace{\Dspace}
\item A \emph{transitive relation.}\vspace{\Dspace}\vspace{\Dspace}
\end{enumerate}
\item Complete the following:

\begin{enumerate}
\item Find all the subsets of $A=\left\{ a,b,c\right\} $.\vspace{\Dspace}\vspace{\Dspace}
\item If $\left|A\right|=4$ how many subsets does $A$ have?\vspace{\Dspace}
\end{enumerate}
\item Let $A=\left\{ 1,2,3\right\} $ and $B=\left\{ 2,4\right\} $.

\begin{enumerate}
\item Find $A\cup B$.\vspace{\Dspace}
\item Find $A\cap B$.\vspace{\Dspace}
\item Find $A-B$.\vspace{\Dspace}
\item Find $A\times B$.\vspace{\Dspace}
\end{enumerate}
\end{enumerate}
\newpage
\begin{enumerate}[resume]
\item Let $A\subseteq B$ and $B\subseteq C$. Prove that $A\subseteq C$.\vfill
\item Prove that if $A\subseteq B$ then $A\cup B=B$.\vfill
\item Prove that $(A\cap B)^{c}\subseteq A^{c}\cup B^{c}$.\vfill
\item Let $f:A\rightarrow B$ be a function from set $A$ to set $B$.

\begin{enumerate}
\item If $\left|A\right|=20$ and $\left|B\right|=21$ can $f$ be one-to-one?
\vspace{\Dspace}
\item If $\left|A\right|=20$ and $\left|B\right|=21$ can $f$ be onto?
\vspace{\Dspace}
\item If $f$ is a constant function and $\left|B\right|>1$, can it be
$1-1$? ($f$ is \emph{constant function }if $f(a)=b$ for all $a\in A$
and some $b\in B$.)\vspace{\Dspace}
\end{enumerate}
\item Complete the following:

\begin{enumerate}[resume]
\item Give an example of a function which is \uline{not} one-to-one.\vspace{\Dspace}
\item Give an example of a function which is \uline{not} onto.\vspace{\Dspace}
\end{enumerate}
\end{enumerate}
\newpage
\begin{enumerate}[resume]
\item Let $f\colon\mathbb{R}\rightarrow\mathbb{R}$ be defined as $f(x)=x^{3}$.

\begin{enumerate}[resume]
\item Prove that $f$ is one-to-one. \vfill
\item Prove that $f$ is onto.\vfill
\end{enumerate}
\item Let $\mathcal{R}=\left\{ \left(a,b\right),\left(b,a\right),\left(b,c\right),\left(a,a\right),\left(b,b\right),\left(c,c\right)\right\} $
be a relation on the set $E=\left\{ a,b,c\right\} $.

\begin{enumerate}[resume]
\item Is $\mathcal{R}$ reflexive? Explain. \vspace{\Dspace}
\item Is $\mathcal{R}$ symmetric? Explain.\vspace{\Dspace}
\end{enumerate}
\item Define the relation $\mathcal{R}$ on the set of even integers as
$x\mathcal{R}y\iff2|x\cdot y$ (i.e., $x\cdot y$ is an even integer).
Show that $\mathcal{R}$ is an equivalence relation.\vfill
\item Complete the following:

\begin{enumerate}
\item Is $16\equiv4\bmod12$?\vspace{\Dspace}
\item Is $5\equiv4\bmod12$?\vspace{\Dspace}
\item Find an integer $x$ such that $0\leq x<4$ and $x\equiv15\bmod4$.\vspace{\Dspace}
\end{enumerate}
\end{enumerate}

\end{document}
